\section{Mathematical formulation}
\label{sec:math}

This section writes down every equation Gate~0 actually solves. Each
subsection starts in ordinary language, then gives the symbols. Quantities
are million tonnes (MMT) per step unless labelled annual; prices are real
dollars per tonne (World Bank Pink Sheet, MUV 2010). Code names match
\texttt{sheaf/dynamic\_crop.py} in the commit that ships with this note.

\subsection{What the model is doing, in one paragraph}

Imagine the world grain market ticking every two weeks. Each country starts
the step with grain in store. Harvest may arrive. People and livestock want
to eat (and, for US maize, ethanol plants want feedstock). Whatever is left
after feeding people and keeping a buffer through the next hungry season can
be offered for export --- unless the government has a ban or tax on the
books. Importers try to buy from their usual suppliers first. Exporters post
a selling price (an ``ask'') and raise it when they sell out, cut it when
grain sits unsold. The world price that we plot is mostly those posted
selling prices, with a small extra kick if the whole system is tighter than
a calm-weather copy of itself. That is the whole spine. Substitution across
grains and governments \emph{choosing} bans are switched off.

\subsection{Clock, countries, and what is in store}
\label{sec:clock}

Time is not annual. There are $T_y=24$ steps per calendar year
($\Delta t\approx 15.2$ days), matching Agrimate~\citep{kuhla2025}. The map
has 17 named countries plus a Rest-of-World bucket. One crop runs at a time:
wheat, then maize, then rice, never all three together here.

At the start of step $t$, country $i$ has stock $S_{i,t}$ (grain already in
silos and pipelines), the world has a price $p_t$, and each exporter has a
posted selling price $q_{i,t}$ (the ask). Harvest $H_{i,t}$ arrives at the
beginning of the step. What the country can use or sell this step is
whatever was in store plus what just came in:
\[
\mathrm{avail}_{i,t}=S_{i,t}+H_{i,t}.
\]

\subsection{Harvest calendars and harvest shocks}
\label{sec:harvest}

\textbf{In plain English.} We do not dump each year's USDA harvest onto the
calendar as a raw pile. First we build a ``normal year'': each country's
typical seasonal harvest, spread over the months when that country actually
cuts grain. Then we ask: was this year high or low \emph{relative to that
country's own slow trend}, not relative to a six-year average that mixes
shocks with the fact that yields drifted up after 2008. The smoother is
LOWESS, a local trend~\citep{cleveland1979}. A year 4\% below trend is a
4\% shock, even if it is 9\% below a boom-inflated sample mean.

Annual USDA PSD production $Y_{i,y}$ is never laid on the calendar as a raw
level. The twin (climatology) path is the score-window \emph{mean} seasonal
harvest: triangular month weights from \texttt{sheaf/seasonal.py} applied to
each country's mean production over the scoring years. Treatment harvest is
that climatology times a per-country LOWESS multiplier,
\begin{equation}
\label{eq:harvest}
H_{i,y}^{\mathrm{ann}}=\overline{H}_i\cdot(1+a_{i,y}),
\qquad
a_{i,y}=\frac{Y_{i,y}-\hat Y_{i,y}}{\hat Y_{i,y}},
\end{equation}
where $\hat Y$ is a local-linear LOWESS trend on a padded PSD history with a
$\sim$10-year window (\texttt{detrend\_anomalies}; tricube weights, local
intercept at the evaluation year). Signed anomalies keep bumper crops as
well as shortfalls (\texttt{shock\_mode=full}).

\textbf{Why not the in-sample mean?} Post-2008 yield growth inflates a
2006--11 average. World wheat 2006 is about $-9\%$ versus that average and
only $-4\%$ versus LOWESS. Using raw levels treats 2006/07 as a crash
(spurious hungry-season spike) and under-weights the 2010 Russian drought,
which sits near a boom-inflated mean. The same bias made 2008 maize look
scarce when it is on-trend. Agrimate uses the same LOWESS idea; we follow it.

Rice calendars are multi-crop (kharif+rabi / early+late) except Vietnam,
whose autumn pulse is kept so the 2008 ban still hits grain that could have
been exported.

\subsection{Demand: food and feed versus ethanol}
\label{sec:demand}

\textbf{In plain English.} Not all use of grain cares about this month's
price. Food and animal feed do: if wheat is expensive, people and feedlots
use a bit less. US maize going into ethanol, in the short run, does not:
blending mandates do not shop for a bargain. We split USDA domestic
consumption into those two piles. Only the food-and-feed pile shrinks when
the world price rises. The ethanol pile is the extra US maize ``food, seed,
and industrial'' use above the pre-mandate 2000--04 average --- a proxy,
because international PSD does not report gallons of ethanol.

PSD domestic use is split
\begin{equation}
C_{i,y}=C^{\mathrm{flex}}_{i,y}+C^{\mathrm{ind}}_{i,y}.
\end{equation}
Flex use (food+feed) faces an isoelastic short-run response: a 1\% price
rise cuts desired use by $|\varepsilon|$ percent. Industrial use is
price-inelastic in the step (a mandate, not a household first-order
condition). For nodes in \texttt{industrial\_nodes} (default: USA maize
only),
\begin{equation}
\label{eq:industrial}
C^{\mathrm{ind}}_{i,y}=\min\Bigl(C_{i,y},\;
\max\bigl(0,\;\mathrm{FSI}_{i,y}-\overline{\mathrm{FSI}}_{i,2000\text{--}04}\bigr)\Bigr).
\end{equation}
That residual is the RFS-era FSI surge~\citep{rfs2005,abbott2011,roberts2013}:
US maize FSI 77~MMT in 2005 to 163~MMT in 2010, while feed fell. Wheat FSI is
flat in this window; rice has no FSI attribute. China's FSI trend is
\emph{not} treated as ethanol.

Official scoring uses the \textbf{average} food-and-feed path over 2006--11,
plus US maize ethanol on (\texttt{use\_demand=False},
\texttt{use\_industrial=True}). Feeding the model each year's actual PSD
food/feed is a sensitivity, not the headline. The calm-weather twin has
average food/feed and no ethanol.

Within-step desired use is
\begin{equation}
\label{eq:desired}
d_{i,t}=C^{\mathrm{flex}}_{i,t}\,(p_t/p_0)^{\varepsilon}+C^{\mathrm{ind}}_{i,t},
\end{equation}
with $\varepsilon<0$ (Table~\ref{tab:params}). $p_0$ is the mean real Pink
Sheet in the reference year (2006 for the 2006--11 hindcast; 2021 for the
Ukraine-war exercise).

\subsection{Hungry-season buffers, then export offers}
\label{sec:lean}

\textbf{In plain English.} Countries do not export every last tonne. They
keep enough grain to eat until the next big harvest arrives, plus a small
standing safety stock. Only the surplus after that buffer is offered to the
world market --- and a ban or tax then shaves that offer. The look-ahead is
crude: blend this year's harvest with a normal seasonal path, count months
until the next global harvest pulse, and hold the gap. This is a stocking
rule of thumb, not a futures market and not a rational-expectations storage
model~\citep{deaton1992}.

Let $h_t$ be steps to the next global harvest pulse (cumulative future world
harvest $\ge 12\%$ of mean annual world $H$, capped at 24 steps). Expected
harvest for that look-ahead blends realized and seasonal paths,
\begin{equation}
H^{\mathrm{exp}}_{i,t}=\phi\,H_{i,t}+(1-\phi)\,H^{\mathrm{seas}}_{i,t}.
\end{equation}
Rolling sums $C^{\mathrm{ahead}}$, $H^{\mathrm{ahead}}$ cover the horizon.
Lean cover (grain needed to bridge to the next pulse) and the stock target
are
\begin{equation}
\label{eq:lean}
L_{i,t}=\max\bigl(0,\,C^{\mathrm{ahead}}_{i,t}+C_{i,t}-H^{\mathrm{ahead}}_{i,t}-H^{\mathrm{exp}}_{i,t}\bigr),
\qquad
T_{i,t}=L_{i,t}+s_i,
\end{equation}
with safety $s_i=\mathtt{stu\_target}\cdot C_i^{\mathrm{ann}}$ (a target
stock-to-use ratio times annual consumption). Purchase demand on the world
market and export offers are
\begin{align}
\label{eq:demand-offers}
D_{i,t}&=\max(0,d_{i,t}-\mathrm{avail}_{i,t})
+\lambda\max\bigl(0,\,T_{i,t}-\max(0,\mathrm{avail}_{i,t}-d_{i,t})\bigr),\\
O_{i,t}&=\max(0,\mathrm{avail}_{i,t}-d_{i,t}-T_{i,t})\,(1-\tau_{i,t}).
\end{align}
In words: if you do not have enough to eat, you try to import the gap, and
you also try to rebuild toward the buffer at speed $\lambda$. If you have
more than food plus the buffer, you offer the extra, then cut that offer by
the export restriction $\tau$ (\S\ref{sec:amis}). $\lambda$ is a
partial-adjustment speed, not a storage Euler equation.

This is Agrimate-aligned commercial cover plus an exogenous restriction.
Continuity with TWIST/Agrimate is the design intent~\citep{schewe2017,kuhla2025};
a departure from Deaton--Laroque is not scored as an error.

\subsection{How much grain a country can hold}
\label{sec:warehouse}

\textbf{In plain English.} Silos are not infinite. Each country has a
ceiling on carry (a maximum stock-to-use) plus, for wheat and maize, room
for about six months of grain in the pipeline --- harvest that has been cut
but is not yet a quiet carry-out stock. If a harvest pulse arrives in one
step, that pulse is allowed to sit under the ceiling so the model does not
pretend the grain vanished on intake day. Rice is different: it is harvested
in several pulses through the year, so we do \emph{not} treat every monsoon
step as extra silo space (that stored the wet season as if it were
warehouses and made December stocks look 60\% too fat). If stocks still
overshoot the ceiling, we drain the overflow slowly rather than deleting it
in one step --- a hard dump was throwing away about a tenth of the wheat and
maize harvest and then leaving importing countries hungry in the spring.

Carry capacity at step $t$ is
\begin{equation}
\label{eq:warehouse}
W_{i,t}=\mathtt{max\_stu}\,C_i^{\mathrm{ann}}
+\mathtt{pipeline\_max\_steps}\cdot C_i^{\mathrm{ann}}/24
+\mathbf{1}_{\{\mathtt{pipeline}>0\}}H_{i,t}.
\end{equation}
Pulse crops (wheat/maize, \texttt{pipeline\_max\_steps=12}) pad $W$ with
same-step harvest. Rice (\texttt{pipeline=0}) clips toward carry only.
Overflow is soft-clipped,
\[
S_{i,t+1}\leftarrow S_{i,t+1}-\mathtt{warehouse\_lambda}\cdot\max(0,S_{i,t+1}-W_{i,t}).
\]
The pipeline term is a \emph{constant} cap (not shrunk as the next harvest
approaches). Seed stocks at spin-up are clipped to
$\max(\mathtt{max\_stu}\,C,\,1.5s)$; the within-step ceiling is
\texttt{max\_stu}\,C, not $1.5s$.

\subsection{Export bans and taxes from the AMIS diary}
\label{sec:amis}

\textbf{In plain English.} We do not estimate who banned exports. We read
the AMIS/OECD policy diary: India banned rice, Russia banned wheat, and so
on~\citep{amis_oecd,kuhla2025}. Those records are labels (prohibition, quota,
tax), not tonnes. We map each label to a fraction of the export \emph{offer}
that is removed --- a ban takes 95\%, a tax 50\%, matching Agrimate. If two
policies overlap, the stricter one wins. Mixed ``maize + wheat'' rows are
dropped rather than guessed. The picture of ``grain locked in the
restricting country'' is then a model object (what would have been offered
without the cut), not FAOSTAT's observed crisis trade.

AMIS/OECD records policy \emph{classes}, not tonnes withheld. Overlapping
policies on the same country take the max cut:
\begin{center}
\small
\begin{tabular}{lc}
\toprule
AMIS \texttt{PolicyMeasure\_Name} & $\tau$ (fraction of offer removed) \\
\midrule
Export prohibition & $0.95$ \\
Export quota & $0.70$ \\
Export tax & $0.50$ \\
Minimum export price & $0.35$ \\
Licensing requirement & $0.25$ \\
Restriction on customs clearance point & $0.15$ \\
\bottomrule
\end{tabular}
\end{center}
This map is structural data, not a free parameter, and is not refit per
crisis. Grain withheld in Fig.~3 analogues is
$O_{i,t}\tau_{i,t}/(1-\tau_{i,t})$.

\subsection{Who trades with whom, and the prices exporters post}
\label{sec:armington}

\textbf{In plain English.} Grain does not splash into a single world pool
and get reallocated to whoever would pay most. Importers have usual
suppliers (Egypt buys a lot of wheat from the Black Sea; Japan from the US
and Australia). That stickiness is an Armington assumption, named for
Armington's 1969 idea that goods from different origins are imperfect
substitutes~\citep{armington1969}. We take the usual supplier mix from
FAOSTAT trade in a quiet window (2006--07 for the food-crisis hindcast;
2019--21 for Ukraine 2022) and \emph{prescribe} it. We do not let a new
network grow out of shipping costs; that is a later paper.

Each fortnight we match export offers to import demand on those usual links,
taking the short side of each pair (you cannot ship more than the exporter
offered or more than the importer wanted from that origin). A small leftover
pool (15\% of unmet demand) can be filled by whoever still has grain --- a
safety valve so a ban does not leave an importer with literally zero
rerouting. Separately, each exporter posts a selling price. If they sold
everything they offered, they raise the ask next step; if grain sat unsold,
they cut it. If a big usual supplier is blocked, surviving exporters mark
prices up a bit (Agrimate's ``the others know they can charge more''
channel). Those posted asks, not a central auctioneer, are most of the
world price in the next subsection.

FAOSTAT E0~\citep{faostat_trade} supplies destination shares $A_{ij}$
(where $i$'s exports go) and source shares $S_{ij}$ (where $j$'s imports
come from), with zero diagonal (no recorded self-trade). FAOSTAT rice E0 in
the crisis files has a \textbf{zero Vietnam row}.
\texttt{load\_trade\_shares} fills the destination mix from 2019--21 and
scales the row to 12\% of crisis-window export total (Vietnam's observed
world-rice-export share). That is a data repair, not a 2008 price knob.

Asks reweight destinations toward cheaper posters,
\begin{equation}
\label{eq:ask-reweight}
\tilde A_{ij}\propto A_{ij}\,(p_0/q_{i,t})^{\gamma}
\quad\text{(rows renormalised)},
\end{equation}
then preferred links match as
\begin{equation}
\label{eq:armington}
\mathrm{ship}_{ij}=\min\bigl(O_{i,t}\tilde A_{ij},\,D_{j,t}S_{ij}\bigr).
\end{equation}
A residual pool may fill at most fraction $\nu=0.15$ of leftover demand,
allocated in proportion to residual offers.

Fill rates update asks (sold-out $\Rightarrow$ raise; leftover $\Rightarrow$
cut), with a survivor markup when preferred sources are blocked:
\begin{equation}
\label{eq:ask}
q_{i,t+1}
=\bigl[(1-\beta)\,q_{i,t}\exp\bigl(\alpha(\mathrm{fill}_{i,t}-\theta)
+\alpha_r b_t\cdot\mathbf{1}_{O_{i,t}>0}\bigr)
+\beta\,p_t\bigr],
\end{equation}
clipped to $[0.45\,p_0,\,2.8\,p_0]$. Fill is shipped/$O$ (target
$\theta=0.70$ if $O=0$). $b_t$ is the preferred-source block fraction
(\S\ref{sec:price}). $\alpha_r=0$ recovers ``only my own fill rate matters.''
Shared $\alpha_r=0.80$ is the smallest value at which turning on maize
export cuts, by themselves, does not \emph{lower} the 14-month mean world
price; it is not fit to 2008 peaks.

Physical update after trade: consumption cannot exceed what you wanted, and
stocks are availability minus exports minus consumption plus imports, then
the silo clip of \S\ref{sec:warehouse}.
\begin{align}
\mathrm{cons}_{i,t}&=\min\bigl(d_{i,t},\,\max(0,\mathrm{avail}_{i,t}-\mathrm{shipped}_i+\mathrm{received}_i)\bigr),\\
S_{i,t+1}&=\max\bigl(0,\mathrm{avail}_{i,t}-\mathrm{shipped}_i-\mathrm{cons}_{i,t}+\mathrm{received}_i\bigr).
\end{align}

\subsection{World price}
\label{sec:price}

\textbf{In plain English.} The number we plot as ``world price'' is not a
magic scarcity formula and not a competitive auction. Most of it is the
tonnage-weighted average of the selling prices exporters actually posted
this step. A smaller part asks: is free grain in the world tighter than in
the calm-weather twin? Free grain means stocks minus what is needed for the
hungry season minus grain trapped behind an export ban (otherwise a ban
would look like a glut sitting in the restricting country). If the treatment
run matches the twin on free grain, unmet import demand, and blocked
suppliers, the target price is exactly the reference $p_0$ --- that is why
the grey baseline in the figures is a flat line, unlike Agrimate's wiggly
unperturbed run. If the world is tighter, the scarcity piece rises; if it is
looser, it falls, so cheap grain can be eaten rather than stored forever.
The plotted price then eases part-way from last step toward that target.
None of those tightness weights is estimated separately on 2008 and 2010.

World-accessible free stocks:
\begin{equation}
\label{eq:free}
\mathrm{locked}_t=\sum_i\tau_{i,t}\max(0,S_{i,t+1}-T_{i,t}),
\quad
\mathrm{free}_t=\sum_i S_{i,t+1}-\sum_i L_{i,t}-\mathrm{locked}_t.
\end{equation}
A path-matched \textbf{twin} is simulated once: seasonal-mean harvest, mean
flex, zero industrial, no export cuts, same starting stocks. Let
$F^{\mathrm{twin}}_t$ and $u^{\mathrm{twin}}_t$ be its free stocks and unmet
fractions. Unmet demand and blockage on the treatment path are
\begin{equation}
u_t=1-\frac{\sum_i\mathrm{received}_{i,t}}{\max(\sum_i D_{i,t},\epsilon)},
\quad
b_t=\frac{\sum_{i,j} S_{ij}\,\tau_{i,t}\,D_{j,t}}{\max(\sum_j D_{j,t},\epsilon)},
\quad
\Delta u_t=\max(0,u_t-u^{\mathrm{twin}}_t).
\end{equation}
Trade-weighted ask $p^{\mathrm{tr}}_t=\sum_i q_{i,t}\mathrm{shipped}_{i,t}/\sum_i\mathrm{shipped}_{i,t}$
(or $p_t$ if no trade). Scarcity
\begin{equation}
\label{eq:scar}
p^{\mathrm{scar}}_t=p_0\cdot r_t^{\eta}
\cdot\bigl(1+\kappa_u\Delta u_t+\kappa_b b_t\bigr),
\quad
r_t=\frac{F^{\mathrm{twin}}_t+f}{\mathrm{free}_t+f},
\end{equation}
with a small positive shift $f$ so the ratio is defined when free stocks
cross zero, and $\eta$ the same in surplus and shortage (muting abundance
stored gluts). Then
\begin{equation}
\label{eq:pstar}
p^\star_t=\omega\,p^{\mathrm{tr}}_t+(1-\omega)\,p^{\mathrm{scar}}_t,
\qquad
p_t=\rho\,p_{t-1}+(1-\rho)\,p^\star_t,
\end{equation}
clipped to $[60,1200]$ \$/t. If free grain, unmet demand, and blockage match
the twin (\textbf{calm}), $p^\star_t=p_0$ by construction
(\texttt{assert\_twin\_identity}).

$\omega$, $\eta$, $\kappa_u$, $\kappa_b$, $\lambda$, $\rho$, and the ask
gains are reduced-form and sign-constrained. They are shared across years
and are \textbf{not} fit crisis-by-crisis.

\subsection{How a run is started}
\label{sec:orch}

\textbf{In plain English.} Before scoring 2006--11 we seed stocks from 2005
USDA ending stocks, then run two quiet years (normal harvest, no bans, no
ethanol) so the silos are not starting from a cold open. We then simulate
the twin once, and the treatment with whichever switches are on. Official
matched scoring turns harvest shocks on, AMIS on, average food/feed, and US
maize ethanol on.

\texttt{run\_crop\_dynamics} (i)~clips seed stocks from PSD year
\texttt{stock\_seed\_year} (2005 for 2006--11; 2020 for 2021--23);
(ii)~spins two climatology years with $\tau=0$ and industrial off;
(iii)~runs the twin; (iv)~runs the treatment with the chosen
\texttt{use\_amis}/\texttt{use\_shocks}/\texttt{use\_demand}/\texttt{use\_industrial}
flags.
